% Options for packages loaded elsewhere
% Options for packages loaded elsewhere
\PassOptionsToPackage{unicode}{hyperref}
\PassOptionsToPackage{hyphens}{url}
\PassOptionsToPackage{dvipsnames,svgnames,x11names}{xcolor}
%
\documentclass[
  letterpaper,
  DIV=11,
  numbers=noendperiod]{scrartcl}
\usepackage{xcolor}
\usepackage{amsmath,amssymb}
\setcounter{secnumdepth}{-\maxdimen} % remove section numbering
\usepackage{iftex}
\ifPDFTeX
  \usepackage[T1]{fontenc}
  \usepackage[utf8]{inputenc}
  \usepackage{textcomp} % provide euro and other symbols
\else % if luatex or xetex
  \usepackage{unicode-math} % this also loads fontspec
  \defaultfontfeatures{Scale=MatchLowercase}
  \defaultfontfeatures[\rmfamily]{Ligatures=TeX,Scale=1}
\fi
\usepackage{lmodern}
\ifPDFTeX\else
  % xetex/luatex font selection
\fi
% Use upquote if available, for straight quotes in verbatim environments
\IfFileExists{upquote.sty}{\usepackage{upquote}}{}
\IfFileExists{microtype.sty}{% use microtype if available
  \usepackage[]{microtype}
  \UseMicrotypeSet[protrusion]{basicmath} % disable protrusion for tt fonts
}{}
\makeatletter
\@ifundefined{KOMAClassName}{% if non-KOMA class
  \IfFileExists{parskip.sty}{%
    \usepackage{parskip}
  }{% else
    \setlength{\parindent}{0pt}
    \setlength{\parskip}{6pt plus 2pt minus 1pt}}
}{% if KOMA class
  \KOMAoptions{parskip=half}}
\makeatother
% Make \paragraph and \subparagraph free-standing
\makeatletter
\ifx\paragraph\undefined\else
  \let\oldparagraph\paragraph
  \renewcommand{\paragraph}{
    \@ifstar
      \xxxParagraphStar
      \xxxParagraphNoStar
  }
  \newcommand{\xxxParagraphStar}[1]{\oldparagraph*{#1}\mbox{}}
  \newcommand{\xxxParagraphNoStar}[1]{\oldparagraph{#1}\mbox{}}
\fi
\ifx\subparagraph\undefined\else
  \let\oldsubparagraph\subparagraph
  \renewcommand{\subparagraph}{
    \@ifstar
      \xxxSubParagraphStar
      \xxxSubParagraphNoStar
  }
  \newcommand{\xxxSubParagraphStar}[1]{\oldsubparagraph*{#1}\mbox{}}
  \newcommand{\xxxSubParagraphNoStar}[1]{\oldsubparagraph{#1}\mbox{}}
\fi
\makeatother


\usepackage{longtable,booktabs,array}
\newcounter{none} % for unnumbered tables
\usepackage{calc} % for calculating minipage widths
% Correct order of tables after \paragraph or \subparagraph
\usepackage{etoolbox}
\makeatletter
\patchcmd\longtable{\par}{\if@noskipsec\mbox{}\fi\par}{}{}
\makeatother
% Allow footnotes in longtable head/foot
\IfFileExists{footnotehyper.sty}{\usepackage{footnotehyper}}{\usepackage{footnote}}
\makesavenoteenv{longtable}
\usepackage{graphicx}
\makeatletter
\newsavebox\pandoc@box
\newcommand*\pandocbounded[1]{% scales image to fit in text height/width
  \sbox\pandoc@box{#1}%
  \Gscale@div\@tempa{\textheight}{\dimexpr\ht\pandoc@box+\dp\pandoc@box\relax}%
  \Gscale@div\@tempb{\linewidth}{\wd\pandoc@box}%
  \ifdim\@tempb\p@<\@tempa\p@\let\@tempa\@tempb\fi% select the smaller of both
  \ifdim\@tempa\p@<\p@\scalebox{\@tempa}{\usebox\pandoc@box}%
  \else\usebox{\pandoc@box}%
  \fi%
}
% Set default figure placement to htbp
\def\fps@figure{htbp}
\makeatother





\setlength{\emergencystretch}{3em} % prevent overfull lines

\providecommand{\tightlist}{%
  \setlength{\itemsep}{0pt}\setlength{\parskip}{0pt}}



 


\KOMAoption{captions}{tableheading}
\makeatletter
\@ifpackageloaded{caption}{}{\usepackage{caption}}
\AtBeginDocument{%
\ifdefined\contentsname
  \renewcommand*\contentsname{Table of contents}
\else
  \newcommand\contentsname{Table of contents}
\fi
\ifdefined\listfigurename
  \renewcommand*\listfigurename{List of Figures}
\else
  \newcommand\listfigurename{List of Figures}
\fi
\ifdefined\listtablename
  \renewcommand*\listtablename{List of Tables}
\else
  \newcommand\listtablename{List of Tables}
\fi
\ifdefined\figurename
  \renewcommand*\figurename{Figure}
\else
  \newcommand\figurename{Figure}
\fi
\ifdefined\tablename
  \renewcommand*\tablename{Table}
\else
  \newcommand\tablename{Table}
\fi
}
\@ifpackageloaded{float}{}{\usepackage{float}}
\floatstyle{ruled}
\@ifundefined{c@chapter}{\newfloat{codelisting}{h}{lop}}{\newfloat{codelisting}{h}{lop}[chapter]}
\floatname{codelisting}{Listing}
\newcommand*\listoflistings{\listof{codelisting}{List of Listings}}
\makeatother
\makeatletter
\makeatother
\makeatletter
\@ifpackageloaded{caption}{}{\usepackage{caption}}
\@ifpackageloaded{subcaption}{}{\usepackage{subcaption}}
\makeatother
\usepackage{bookmark}
\IfFileExists{xurl.sty}{\usepackage{xurl}}{} % add URL line breaks if available
\urlstyle{same}
\hypersetup{
  pdftitle={01\_midterm\_2026},
  colorlinks=true,
  linkcolor={blue},
  filecolor={Maroon},
  citecolor={Blue},
  urlcolor={Blue},
  pdfcreator={LaTeX via pandoc}}


\title{01\_midterm\_2026}
\author{}
\date{}
\begin{document}
\maketitle


\begin{enumerate}
\def\labelenumi{\arabic{enumi}.}
\tightlist
\item
  A company is developing a new product. Let \(X_n\) denote the
  development score after period \(n\). At each period, the score
  increases by one with probability \(0.5\) and decreases by one with
  probability \(0.5\).
\end{enumerate}

The project is \textbf{successfully completed} when the score reaches
\(10\) and \textbf{abandoned} when the score reaches \(0\). Suppose that
the initial development score is \(X_0=4\).

Let

\[
T=\inf\{n\geq0:X_n\in\{0,10\}\}
\]

be the number of periods until the project is completed or abandoned.

Using the results for a simple random walk with absorbing barriers at
\(0\) and \(N\):

\begin{enumerate}
\def\labelenumi{\arabic{enumi}.}
\item
  Calculate the probability that the project is abandoned before it is
  completed.
\item
  Calculate the expected number of periods until the project is
  completed or abandoned.
\end{enumerate}

\begin{center}\rule{0.5\linewidth}{0.5pt}\end{center}

\begin{enumerate}
\def\labelenumi{\arabic{enumi}.}
\setcounter{enumi}{1}
\tightlist
\item
  A motor insurer operates a No Claims Discount (NCD) system with three
  levels: \(0%
  \), \(15%
  \), and \(40%
  \).
\end{enumerate}

The policyholder's discount level changes according to the number of
claims made during the year:

\begin{itemize}
\tightlist
\item
  \textbf{No claims:} Move up one level of discount, or remain at the
  maximum level.
\item
  \textbf{Exactly one claim:} Move down one level, or remain at the
  minimum level.
\item
  \textbf{Two or more claims:} Move to, or remain at, the minimum level.
\end{itemize}

Let

\[
a=P(\text{exactly one claim in a year})
\]

and

\[
b=P(\text{two or more claims in a year}).
\]

It is known that the probability of making two or more claims is \(25%
\) of the probability of making exactly one claim.

\begin{enumerate}
\def\labelenumi{\arabic{enumi}.}
\item
  Express \(b\) and the probability of making no claim in terms of
  \(a\).
\item
  Taking the states to be \(0%
  \), \(15%
  \), and \(40%
  \), write down the transition matrix \(P\) in terms of \(a\). If
  \(\boldsymbol{\pi}\) denotes the long-run probability vector, then it
  satisfies

  \[
  \boldsymbol{\pi}P=\boldsymbol{\pi}.
  \]

  It is known that the long-run probability of being at the \(15%
  \) level is equal to the long-run probability of being at the \(40%
  \) level. Use this information to show that

  \[
  a=0.4.
  \]
\item
  Using the value of \(a\) found in Question 2, write down the
  transition matrix \(P\). Calculate the probability that a policyholder
  who is at the \(0%
  \) level this year will be at the \(40%
  \) level after three years.
\item
  Calculate the probability that a policyholder who is at the \(0%
  \) level this year will be at the \(40%
  \) level after 10 years. Briefly explain why this probability is close
  to the long-run probability of being at the \(40%
  \) level.
\end{enumerate}

\begin{center}\rule{0.5\linewidth}{0.5pt}\end{center}

\begin{enumerate}
\def\labelenumi{\arabic{enumi}.}
\setcounter{enumi}{2}
\tightlist
\item
  A financial institution classifies a company according to its credit
  rating. The possible states of a Markov chain are
\end{enumerate}

\[
\{A,B,C,D,E\},
\]

where state \(E\) represents default.

The transition matrix is

\[
P=
\begin{pmatrix}
0.6 & 0.3 & 0.1 & 0 & 0\\
0.2 & 0.6 & 0.2 & 0 & 0\\
0 & 0 & 0.7 & 0.3 & 0\\
0 & 0 & 0.4 & 0.6 & 0\\
0 & 0 & 0 & 0 & 1
\end{pmatrix}.
\]

\begin{enumerate}
\def\labelenumi{\arabic{enumi}.}
\item
  Suppose a company is currently in state \(A\). Calculate:

  \begin{enumerate}
  \def\labelenumii{\alph{enumii}.}
  \item
    the probability that its credit rating will be in state \(C\) after
    two years;
  \item
    the probability that it will have defaulted after three years.
  \end{enumerate}

  \emph{Hint:} For part (a), consider all possible paths from \(A\) to
  \(C\) involving exactly two transitions. For part (b), first identify
  how state \(E\) can be reached.
\item
  Determine which pairs of states communicate with each other. Give a
  brief justification for your answer.
\item
  Classify the state space into its communication classes. Identify
  which communication classes are closed.
\item
  State which states are transient and which are recurrent.
\end{enumerate}

\begin{center}\rule{0.5\linewidth}{0.5pt}\end{center}

\begin{enumerate}
\def\labelenumi{\arabic{enumi}.}
\item
  \textbf{Deterministic and stochastic processes}

  State whether each of the following is deterministic or stochastic,
  and briefly justify your answer:

  \begin{enumerate}
  \def\labelenumii{\alph{enumii}.}
  \item
    The value of a savings account that earns a fixed interest rate of
    \(5%
    \) per year with no deposits or withdrawals.
  \item
    The number of claims received by an insurer during a year.
  \item
    The temperature recorded at a particular location at noon each day.
  \end{enumerate}
\item
  \textbf{Time-homogeneous and time-inhomogeneous Markov chains}

  A Markov chain has transition probabilities satisfying

  \[
  P(X_{n+1}=j\mid X_n=i)=p_{ij}^{(n)}.
  \]

  State the condition on \(p_{ij}^{(n)}\) under which the Markov chain
  is time-homogeneous. Give one example of a situation in which the
  transition probabilities may depend on time.
\item
  \textbf{Periodicity and aperiodicity}

  Consider a Markov chain with transition matrix
\end{enumerate}

\[
  P=
  \begin{pmatrix}
  0 & 0 & 1\\
  1 & 0 & 0\\
  0 & 1 & 0
  \end{pmatrix}
  \]

\begin{enumerate}
\def\labelenumi{\alph{enumi}.}
\item
  Determine the period of each state.
\item
  State whether the Markov chain is periodic or aperiodic.
\item
  Does the chain have a stationary distribution? If so, find it.
\item
  Starting from state \(1\), explain why the state probabilities do not
  converge to the stationary distribution.
\end{enumerate}




\end{document}
